% BTS CIEL 1 — Réseaux — 04 — ARP
% Compiler deux fois avec XeLaTeX depuis fichiers/.
\documentclass[aspectratio=169,11pt]{beamer}
\usepackage{fontspec}
\setmainfont{Latin Modern Sans}
\setsansfont{Latin Modern Sans}
\setmonofont{Latin Modern Mono}
\usepackage{polyglossia}
\setmainlanguage{french}
\usepackage{tikz}
\usetikzlibrary{arrows.meta,positioning,calc,shapes.geometric}
\usepackage{listings}
\usepackage{xcolor}
\definecolor{foret}{HTML}{1C322D}
\definecolor{creme}{HTML}{F8F3EE}
\definecolor{or}{HTML}{EBB552}
\definecolor{rose}{HTML}{F1CDBE}
\definecolor{vertclair}{HTML}{DDE9DF}
\definecolor{encre}{HTML}{24372E}
\setbeamercolor{normal text}{fg=creme,bg=foret}
\setbeamercolor{frametitle}{fg=or,bg=foret}
\setbeamercolor{title}{fg=creme}
\setbeamercolor{block title}{fg=foret,bg=or}
\setbeamercolor{block body}{fg=encre,bg=creme}
\setbeamercolor{alerted text}{fg=or}
\setbeamercolor{structure}{fg=or}
\setbeamertemplate{navigation symbols}{}
\setbeamertemplate{footline}{\hfill\color{or}\insertframenumber\hspace{0.5cm}\vspace{0.25cm}}
\setbeamertemplate{frametitle}{\vspace{0.35cm}\insertframetitle\par\vspace{0.15cm}\color{or}\rule{\linewidth}{0.6pt}}
\setlength{\parskip}{5pt}
\lstdefinestyle{webcode}{basicstyle=\ttfamily\small\color{encre},keywordstyle=\color{foret},stringstyle=\color{foret},
  backgroundcolor=\color{creme},frame=single,rulecolor=\color{or},framerule=1pt,
  xleftmargin=5pt,xrightmargin=5pt,columns=fullflexible,keepspaces=true,showstringspaces=false}
\newcommand{\question}[1]{\begin{beamercolorbox}[rounded=true,sep=12pt]{block body}\Large\textbf{#1}\end{beamercolorbox}}
\newcommand{\etiquette}[1]{{\color{or}\small\MakeUppercase{#1}}}
% Les schémas ont une surface claire, également lisible en noir et blanc.
\tikzset{
  boite/.style={draw=encre,fill=creme,text=encre,rounded corners,
    align=center,text width=2.8cm,minimum height=1cm},
  flux/.style={-{Stealth},thick,draw=encre},
  every node/.style={text=encre,font=\small}}
\newenvironment{schema}{%
  \begin{beamercolorbox}[sep=8pt,rounded=true]{block body}
  \centering\begin{tikzpicture}[>=Stealth]
}{\end{tikzpicture}\end{beamercolorbox}}
\title{Le protocole ARP}
\subtitle{BTS CIEL 1 · Réseaux · 04 — ARP}
\author{}
\date{}
\begin{document}
\begin{frame}[plain]
\vfill\etiquette{Trouver une adresse sur le réseau local}\par\bigskip
{\Huge\bfseries Le protocole ARP}\par\bigskip
{\Large BTS CIEL 1 · Réseaux · 04 — ARP}\par\bigskip
\begin{schema}
\node[boite,text width=4cm,minimum height=2cm,font=\large] (a) {Adresse IPv4};
\node[boite,text width=4cm,minimum height=2cm,font=\large,right=2cm of a] (b) {Adresse MAC};
\draw[flux] (a)--node[above,font=\normalsize]{ARP} (b);
\end{schema}\vfill
\end{frame}
\begin{frame}{Trouver une collègue}
\question{Comment trouver une collègue dans un open space ?}
\begin{schema}
\begin{scope}[scale=0.9,transform shape]
\node[boite,text width=2cm,minimum height=0.9cm] (a) at (0,0) {Vous};
\node[boite,text width=2.5cm] (b) at (8,1.5) {Collègue 1};
\node[boite,text width=2.5cm] (c) at (8,0) {Collègue recherchée};
\node[boite,text width=2.5cm] (d) at (8,-1.5) {Collègue 3};
\draw[flux,densely dotted] (a.north east)--(b.west);
\draw[flux,densely dotted] ([yshift=3mm]a.east)--([yshift=3mm]c.west);
\draw[flux,densely dotted] (a.south east)--(d.west);
\draw[flux,dashed] ([yshift=-3mm]c.west)--([yshift=-3mm]a.east);
\node[font=\small,fill=creme] at (4,1.35) {1. Question à voix haute};
\node[font=\small,fill=creme] at (4,-0.55) {2. Réponse};
\end{scope}
\end{schema}
Sans réponse, vous ne savez pas où la trouver.
\end{frame}
\begin{frame}{Une adresse manque}
\question{Comment envoyer une trame à PC B ?}
\begin{schema}
\node[boite,text width=3.7cm,minimum height=2.7cm,font=\normalsize] (a) at (0,0) {PC A\\\texttt{192.168.10.11/24}};
\node[boite,text width=3.7cm,minimum height=2.7cm,font=\normalsize] (b) at (9,0) {PC B\\\texttt{192.168.10.12/24}};
\draw[flux] (a)--node[above,align=center]{IPv4 connue\\MAC inconnue} (b);
\end{schema}
La trame Ethernet a besoin d'une adresse MAC de destination.
\end{frame}
\begin{frame}{Une recherche sur le réseau local}
\begin{schema}
\node[boite,text width=3.3cm] (a) at (0,0) {PC A\\\texttt{02:00:00:00:00:A1}};
\node[boite,text width=2.2cm] (s) at (4,0) {Commutateur};
\node[boite,text width=2.8cm] (b) at (9,0.9) {PC B\\\texttt{192.168.10.12}};
\node[boite,text width=2cm] (c) at (9,-0.9) {PC D\\ignore la requête};
\draw[flux,densely dotted] ([yshift=3mm]a.east)--([yshift=3mm]s.west);
\draw[flux,densely dotted] ([yshift=3mm]s.east)--([yshift=3mm]b.west);
\draw[flux,densely dotted] ([yshift=-3mm]s.east)--(c.west);
\draw[flux,dashed] ([yshift=-3mm]b.west)--([yshift=1mm]s.east);
\draw[flux,dashed] ([yshift=-3mm]s.west)--([yshift=-3mm]a.east);
\node[align=center,font=\small] at (4,1.65) {1. Requête ARP à la MAC de diffusion \texttt{FF:FF:FF:FF:FF:FF} :\\qui a \texttt{192.168.10.12} ?};
\node[align=center,font=\small] at (4,-1.25) {2. Réponse ARP (unicast) :\\\texttt{192.168.10.12} est à \texttt{02:00:00:00:00:3C}};
\end{schema}
{\small La requête contient l'IPv4 et la MAC de PC A : PC B peut lui répondre.\\
Pointillés : requête ; tirets : réponse. Unicast : un seul destinataire.\\
Même MAC de diffusion : \texttt{FF:FF:FF:FF:FF:FF}, \texttt{ff-ff-ff-ff-ff-ff} (Windows),\\
\texttt{FFFF.FFFF.FFFF} (Packet Tracer).}
\end{frame}
\begin{frame}{Le cache évite de redemander}
\begin{schema}
\begin{scope}[scale=1.1,transform shape]
\node[boite,text width=4.5cm,minimum height=2.8cm,font=\normalsize] (a) at (0,0) {Cache de PC A\\\texttt{192.168.10.12}\\\texttt{02:00:00:00:00:3C}};
\node[boite,text width=4.5cm,minimum height=2.8cm,font=\normalsize] (b) at (7.3,0) {Trame suivante\\MAC de destination :\\\texttt{02:00:00:00:00:3C}};
\draw[flux] (a)--node[above]{1. réutiliser} (b);
\end{scope}
\end{schema}
Le cache ARP (ou table ARP) conserve les correspondances IPv4–MAC.\\
Les entrées dynamiques sont apprises automatiquement et s'effacent après un délai. Sous Windows, consultez-les avec \texttt{arp -a}.
\end{frame}
\begin{frame}{Utiliser ARP}
\begin{schema}
\begin{scope}[scale=0.95,transform shape]
\node[boite,text width=3.6cm,minimum height=2cm,font=\normalsize] (a) at (0,0) {IPv4 du destinataire\\connue};
\node[boite,text width=3.6cm,minimum height=2cm,font=\normalsize] (b) at (5.3,0) {Requête puis réponse};
\node[boite,text width=3.6cm,minimum height=2cm,font=\normalsize] (c) at (10.6,0) {MAC du destinataire\\connue};
\draw[flux] (a)--(b); \draw[flux] (b)--(c);
\end{scope}
\end{schema}
\begin{block}{Méthode}
1. Cherchez l'IPv4 dans le cache.\\
2. Si elle manque, émettez une requête.\\
3. Lisez la MAC et enregistrez la correspondance dans le cache.\\
4. Utilisez cette MAC pour la trame.
\end{block}
\end{frame}
\begin{frame}{Une destination ailleurs}
\question{Quelle MAC utiliser pour un autre réseau ?}
\begin{schema}
\node[boite,text width=3.6cm,minimum height=2.3cm,font=\normalsize] (a) at (0,0) {PC A\\\texttt{192.168.10.11/24}};
\node[boite,text width=3.1cm,minimum height=2.3cm,font=\normalsize] (g) at (5,0) {Passerelle\\\texttt{192.168.10.1}};
\node[boite,text width=3.1cm,minimum height=2.3cm,font=\normalsize] (d) at (10,0) {Serveur distant\\\texttt{198.51.100.20}};
\draw[flux] (a)--(g); \draw[flux] (g)--(d);
\node[font=\small,align=center] at (2.5,-1.9) {MAC dest. : passerelle\\IPv4 dest. : \texttt{198.51.100.20}};
\end{schema}
PC A cherche la MAC de sa passerelle, pas celle du serveur.
\end{frame}
\begin{frame}{Choisir l'adresse à résoudre}
\begin{schema}
\node[boite,text width=4cm,minimum height=2cm,font=\normalsize] (a) at (0,0) {Même réseau ?\\MAC du destinataire};
\node[boite,text width=4cm,minimum height=2cm,font=\normalsize] (b) at (8,0) {Autre réseau ?\\MAC de la passerelle};
\end{schema}
\begin{block}{Méthode}
1. Appliquez votre masque à votre IPv4 et à l'IPv4 de destination,\\
puis comparez les deux réseaux obtenus.\\
2. Choisissez la destination locale ou la passerelle.\\
3. Consultez le cache ; au besoin, résolvez cette IPv4 avec ARP.
\end{block}
\end{frame}
\begin{frame}{À vous : choisir une adresse}
\question{Quelle IPv4 PC C doit-il résoudre ? Justifiez.}
\begin{schema}
\node[boite,text width=4cm,minimum height=2.6cm,font=\normalsize] (a) at (0,0) {PC C\\\texttt{192.168.20.4/24}\\passerelle : \texttt{192.168.20.1}};
\node[boite,text width=4cm,minimum height=2.6cm,font=\normalsize] (b) at (9,0) {Imprimante\\\texttt{192.168.20.30/24}};
\draw[flux] (a)--node[above]{Document à envoyer} (b);
\end{schema}
\end{frame}
\begin{frame}{Correction : destination locale}
\begin{schema}
\begin{scope}[scale=1.15,transform shape]
\node[boite,text width=4.8cm,minimum height=2.5cm] (a) at (0,0) {Masque : \texttt{255.255.255.0}\\Dernier octet : 4 ET 0 = 0\\30 ET 0 = 0};
\node[boite,text width=4.2cm,minimum height=2.5cm] (b) at (6.8,0) {Même réseau :\\\texttt{192.168.20.0/24}\\Résoudre \texttt{192.168.20.30}};
\draw[flux] (a)--(b);
\end{scope}
\end{schema}
La passerelle ne sert pas pour cette destination locale.
\end{frame}
\begin{frame}{À vous : une destination à atteindre}
\question{Quelle IPv4 PC C résout-il avec ARP ? Justifiez.}
\begin{schema}
\begin{scope}[scale=1.1,transform shape]
\node[boite,text width=4.3cm,minimum height=2.5cm] (a) at (0,0) {PC C\\\texttt{192.168.20.4/25}\\Passerelle : \texttt{192.168.20.1}};
\node[boite,text width=4cm,minimum height=2.5cm] (b) at (7,0) {Destination\\\texttt{192.168.20.130}};
\draw[flux] (a)--(b);
\end{scope}
\end{schema}
\end{frame}
\begin{frame}{Correction : choisir la passerelle}
\begin{schema}
\begin{scope}[scale=1.05,transform shape]
\node[boite,text width=6.1cm,minimum height=2.7cm] (a) at (0,0) {Masque : \texttt{255.255.255.128}\\Dernier octet :\\\texttt{00000100 ET 10000000 = 00000000}\\\texttt{10000010 ET 10000000 = 10000000}};
\node[boite,text width=4.5cm,minimum height=2.7cm] (b) at (7.1,0) {Réseaux différents\\\texttt{192.168.20.0}\\\texttt{192.168.20.128}};
\draw[flux] (a)--(b);
\end{scope}
\end{schema}
PC C résout \texttt{192.168.20.1}, l'adresse de sa passerelle.
\end{frame}
\begin{frame}{À retenir}
\begin{schema}
\node[boite,text width=3.6cm,minimum height=1.8cm,font=\normalsize] (a) {IPv4 du destinataire};
\node[boite,text width=3.6cm,minimum height=1.8cm,font=\normalsize,right=1cm of a] (b) {ARP ou cache};
\node[boite,text width=3.6cm,minimum height=1.8cm,font=\normalsize,right=1cm of b] (c) {MAC du destinataire};
\draw[flux] (a)--(b); \draw[flux] (b)--(c);
\end{schema}
\begin{block}{Quatre repères}
ARP relie l'IPv4 d'un voisin local à sa MAC.\\
La requête est diffusée ; la réponse est envoyée au demandeur.\\
Pour un réseau distant, utilisez la MAC de la passerelle.\\
ARP n'authentifie pas les réponses : une réponse peut être falsifiée.
\end{block}
\end{frame}
\begin{frame}{Sources}
\begin{beamercolorbox}[sep=14pt,rounded=true]{block body}
\textbf{RFC 826} — résolution d'adresses Ethernet\\
\texttt{https://www.rfc-editor.org/rfc/rfc826}\\[8pt]
\textbf{RFC 1122, § 2.3.2.1} — expiration des entrées ARP\\
\texttt{https://www.rfc-editor.org/rfc/rfc1122\#section-2.3.2.1}\\[8pt]
\textbf{Microsoft Learn} — commande Windows \texttt{arp}\\
\scriptsize\texttt{https://learn.microsoft.com/windows-server/administration/windows-commands/arp}
\end{beamercolorbox}
\end{frame}
\end{document}
